% mono/preamble/setup.tex
%
% Base preamble for the Agentic Systems Framework monograph.
% Loaded by main.tex after \documentclass{kaobook}.
%
% Conventions:
% — fontmethod=modern (unicode-math) so we can use fontspec to swap in
%   open-source Garamond + Roboto for the Principia-flavored aesthetic.
% — Math-compat macros for GFM/Obsidian idioms (\lt, \gt, etc.) are defined
%   here rather than reverse-transformed in source: the segment markdown is
%   the source of truth and stays viewer-friendly in both worlds.

%----------------------------------------------------------------------------
% Language & quoting
%----------------------------------------------------------------------------
\usepackage[english]{babel}
\usepackage[english=american]{csquotes}

%----------------------------------------------------------------------------
% Math
%----------------------------------------------------------------------------
% kaobook's modern fontmethod pulls in unicode-math + Libertinus by default;
% we override the math font below with STIX Two Math for the Garamond pairing.
\usepackage{amsmath}
\usepackage{amssymb}
\usepackage{mathtools}
% amsthm-style theorem environments are provided by kaotheorems (loaded in main.tex).

%----------------------------------------------------------------------------
% Fonts — EB Garamond + Roboto Light + STIX Two Math
%
% These names match what fontspec finds via the OS fontconfig database; on
% TeX Live 2025 the EB Garamond and Roboto packages also install the same
% OTFs into the tex tree, so the lookup works on machines without the system
% fonts. STIX Two Math ships with macOS (we use the system copy) and is also
% available via the stix2-otf TL package on Linux.
%----------------------------------------------------------------------------
% Italic text gets a muted teal tint so emphasis and term-introduction
% reads as a quiet visual register, distinct from body-bold.
\setmainfont{EB Garamond}[
  Ligatures         = {TeX, Common},
  Numbers           = OldStyle,
  SmallCapsFeatures = {LetterSpace=4},
  ItalicFeatures    = {Color=5A787C},
]
\setsansfont{Roboto Light}[
  Ligatures   = TeX,
  Scale       = MatchLowercase,
  BoldFont    = {Roboto Medium},
  ItalicFont  = {Roboto Light Italic},
]
% Monospace: Fira Mono — narrower and cleaner than the Libertinus Mono
% kaobook defaults to; reads well for inline code spans and slug names
% (`#some-slug-name`) that appear frequently in segment prose.
%
% Scale and negative LetterSpace tighten the spread (Fira Mono is wide
% by default at MatchLowercase). Color is a muted olive — the same
% "quiet typographic register" approach as italics; mono content reads
% as a sibling channel rather than competing with body bold.
\setmonofont{Fira Mono}[
  Ligatures   = TeX,
  Scale       = 0.85,
  LetterSpace = -4,
  Color       = 6B7660,
]

\setmathfont{STIX Two Math}

%----------------------------------------------------------------------------
% GFM/Obsidian math-compat macros
%
% Source segments use \lt and \gt because raw < and > break GitHub's math
% renderer (interpreted as HTML tags). LaTeX math knows neither command, so
% define them as the obvious operators. \ast is standard LaTeX math — no-op.
%----------------------------------------------------------------------------
\providecommand{\lt}{<}
\providecommand{\gt}{>}

%----------------------------------------------------------------------------
% Microtypography
%----------------------------------------------------------------------------
\usepackage{microtype}

%----------------------------------------------------------------------------
% String predicates — used by segment-environment and status-badge logic.
%----------------------------------------------------------------------------
\usepackage{xstring}

%----------------------------------------------------------------------------
% wasysym — provides \CIRCLE and \Circle (size-paired filled and empty
% circle glyphs) for the segment-stage indicator. Earlier attempts to use
% raw Unicode ●/○ ran into EB Garamond's missing-glyph fallback and a
% font-subset failure on the system STIX Two Text.
%----------------------------------------------------------------------------
\usepackage{wasysym}

%----------------------------------------------------------------------------
% changepage — adjustwidth* environment with signed margin offsets, used
% by \segmentepigraph to extend into the margin column when centering on
% \segmentrulewidth (positive arg indents inward; negative extends past).
%----------------------------------------------------------------------------
\usepackage{changepage}

%----------------------------------------------------------------------------
% pdfpages — \includepdf places a PDF page (or pages) into the document.
% Used to render per-component cover pages from SVG (the build converts
% the SVG to PDF via rsvg-convert before compile).
%----------------------------------------------------------------------------
\usepackage{pdfpages}

%----------------------------------------------------------------------------
% Vertical-space reservation — \needspace{N\baselineskip} prevents a
% segment header from being orphaned at page bottom by forcing a page
% break first when the requested space isn't available.
%----------------------------------------------------------------------------
\usepackage{needspace}

%----------------------------------------------------------------------------
% Tables — tabularx + X columns wrap content to fit \linewidth, so cells
% with long prose break naturally instead of overflowing the page. We
% deliberately avoid xltabular (longtable-based, supports page breaks)
% because its grouping collides with kaobook's margin-note machinery
% (\cmrsideswitch from booktabs becomes undefined at \end{xltabular}
% inside kaobook's pagestyle). Tables therefore don't break across pages;
% any genuinely tall table needs source-side splitting.
% booktabs gives us the Tufte-style rule register; array gives the
% column-modifier prefixes for raggedright/centering/raggedleft X.
% caption + \captionof let us emit a numbered caption above each table
% without forcing it into a `table` float (which would re-flow).
%----------------------------------------------------------------------------
\usepackage{tabularx}
\usepackage{booktabs}
\usepackage{array}      % for >{\raggedright\arraybackslash}X column types
\usepackage{caption}

%----------------------------------------------------------------------------
% Equation + table numbering — both reset per chapter and display with
% roman chapter numerals to match the segment numbering register (e.g.,
% Definition I.4 followed by equation (I.5)).
%----------------------------------------------------------------------------
\numberwithin{equation}{chapter}
\numberwithin{table}{chapter}
\renewcommand{\theequation}{\Roman{chapter}.\arabic{equation}}
\renewcommand{\thetable}{\Roman{chapter}.\arabic{table}}

%----------------------------------------------------------------------------
% Color palette
%
% Cream paper background is the Principia signature — commented out by
% default so v0.1.0 ships on white, which is the preprint convention. Flip
% the comment to bring the warmth back.
%----------------------------------------------------------------------------
\usepackage{xcolor}
\definecolor{paper-cream}    {HTML}{FFFEF2}  % Principia background
\definecolor{accent-deep}    {HTML}{8B0000}  % accent for section heads, refs
\definecolor{rule-gray}      {HTML}{CCCCCC}
\definecolor{status-exact}   {HTML}{2E5A88}  % blue-ish; settled
\definecolor{status-soft}    {HTML}{8B6914}  % muted gold; heuristic/qualitative
\definecolor{status-tentative}{HTML}{8B4A2F} % muted rust; sketch/discussion-grade
\definecolor{link-navy}      {HTML}{1F3A6F}  % cross-refs, citations, URLs

% \pagecolor{paper-cream}  % uncomment for cream paper

%----------------------------------------------------------------------------
% Hyperref color discipline — all cross-refs, citations, and URLs read
% in navy blue. Applied via \hypersetup at \AtBeginDocument so it wins
% over kaobook's own hypersetup calls (kaorefs/kaobiblio set their own
% colors during their own \AtBeginDocument hooks).
%----------------------------------------------------------------------------
\AtBeginDocument{%
  \hypersetup{
    colorlinks = true,
    linkcolor  = link-navy,
    citecolor  = link-navy,
    urlcolor   = link-navy,
    filecolor  = link-navy,
  }%
}

%----------------------------------------------------------------------------
% That's it for the base preamble. Segment environments, status badges, and
% equation-tag margin commands land in their own preamble files in task 2.
%----------------------------------------------------------------------------
